\section{An explicit mixed-color identity at every even weight}
\label{newmixed:sec:identity}

The polynomial reduction also gives a concrete infinite family of
identities.  The following proof stays entirely within three convergent
color families, so it needs no auxiliary divergent coordinate.
Write
\[
 u_{a,b}=\Im\operatorname{Li}_{a,b}(i,-1),\qquad
 \beta(s)=\Im\operatorname{Li}_s(i),\qquad
 \eta_{\mathrm{alt}}(s)=-\operatorname{Li}_s(-1).
\]
Thus $\eta_{\mathrm{alt}}(1)=\log2$, and
$\eta_{\mathrm{alt}}(s)=(1-2^{1-s})\zeta(s)$ for integers $s\ge2$.

\begin{theorem}[One trailing index in the mixed Gaussian family]
\label{newmixed:thm:family}
For every even integer $w\ge2$,
\begin{equation}\label{newmixed:eq:family}
 \boxed{\displaystyle
 u_{w-1,1}
 =\frac w2\,\beta(w)-\log2\,\beta(w-1)
  -\sum_{\substack{1\le a<w\\a\ {
m odd}}}
       \eta_{\mathrm{alt}}(a)\beta(w-a)
  +2^{1-w}\beta(1)\eta_{\mathrm{alt}}(w-1).}
\end{equation}
In particular,
\begin{align}
 u_{1,1}&=\beta(2)-\frac{3\pi}{8}\log2,\label{newmixed:eq:w2}\\
 u_{3,1}&=2\beta(4)-\frac{\pi^3}{16}\log2
                    -\frac{21\pi}{128}\zeta(3),\label{newmixed:eq:w4}\\
 u_{5,1}&=3\beta(6)-\frac{5\pi^5}{768}\log2
                    -\frac{3\pi^3}{128}\zeta(3)
                    -\frac{465\pi}{2048}\zeta(5).
 \label{newmixed:eq:w6}
\end{align}
\end{theorem}
\begin{proof}
Put $n=w-2$ and define degree-$n$ homogeneous polynomials
\begin{align*}
 C(X,Y)&=\sum_{a+b=w}\Im\operatorname{Li}_{a,b}(i,i)X^{a-1}Y^{b-1},\\
 D(X,Y)&=\sum_{a+b=w}u_{a,b}X^{a-1}Y^{b-1},\\
 F(X,Y)&=\sum_{a+b=w}\Im\operatorname{Li}_{a,b}(-1,i)X^{a-1}Y^{b-1},\\
 P(X,Y)&=\sum_{a+b=w}\Im\bigl(\operatorname{Li}_a(-1)
                         \operatorname{Li}_b(i)\bigr)X^{a-1}Y^{b-1},\\
 Q(X,Y)&=\sum_{a+b=w}\Im\bigl(\operatorname{Li}_a(i)
                         \operatorname{Li}_b(i)\bigr)X^{a-1}Y^{b-1},\\
 R(X,Y)&=\sum_{a+b=w}\Im\bigl(\operatorname{Li}_a(i)
                         \operatorname{Li}_b(-i)\bigr)X^{a-1}Y^{b-1},\\
 H(X,Y)&=\sum_{j=0}^{n}X^jY^{n-j}.
\end{align*}
All indices $a,b$ in these sums are positive.  Every single factor is
finite, including those of index one.  The colored stuffle and shuffle
identities give exactly
\begin{align}
 F(X,Y)+D(Y,X)&=P(X,Y)+\beta(w)H(X,Y),\label{newmixed:eq:SFD}\\
 C(X,Y)+C(Y,X)&=Q(X,Y),\label{newmixed:eq:SCC}\\
 C(X+Y,X)-F(X+Y,Y)&=P(X,Y),\label{newmixed:eq:HCF}\\
 D(X+Y,Y)-D(X+Y,X)&=R(X,Y).\label{newmixed:eq:HDD}
\end{align}
The signs in the first and third equations use respectively
$\Im\operatorname{Li}_w(-i)=-\beta(w)$ and
$\Im\operatorname{Li}_{a,b}(-1,-i)
=-\Im\operatorname{Li}_{a,b}(-1,i)$.

Let $d=D(1,0)=u_{w-1,1}$.  Equation~\eqref{newmixed:eq:HDD} at
$(X,Y)=(0,1)$ says $D(1,1)=d+R(0,1)$.  Using
\eqref{newmixed:eq:HCF} at $(0,1)$ and
\eqref{newmixed:eq:SFD} at $(1,1)$ gives
\[
 C(1,0)=-D(1,1)+P(1,1)+P(0,1)+(w-1)\beta(w).
\]
Since $n$ is even, $F(0,-1)=F(0,1)$.  The same equations at
$(1,-1)$ and $(0,1)$ give
\[
 C(0,1)=-d+P(0,1)+P(1,-1)+\beta(w).
\]
Adding and using \eqref{newmixed:eq:SCC} yields the exact identity
\begin{equation}\label{newmixed:eq:two-d}
 2d=P(1,1)+P(1,-1)+2P(0,1)+w\beta(w)-Q(1,0)-R(0,1).
\end{equation}
Because $a+b=w$ is even, $a$ and $b$ have the same parity.  Therefore
\[
 P(1,1)+P(1,-1)
 =-2\sum_{\substack{1\le a<w\\a\ {
m odd}}}\eta_{\mathrm{alt}}(a)\beta(w-a),
 \qquad P(0,1)=-\log2\,\beta(w-1).
\]
Finally,
\begin{align*}
 Q(1,0)+R(0,1)
 &=\Im\left[\operatorname{Li}_1(i)
       \bigl(\operatorname{Li}_{w-1}(i)
             +\operatorname{Li}_{w-1}(-i)\bigr)\right]\\
 &=2\beta(1)\Re\operatorname{Li}_{w-1}(i)
 =-2^{2-w}\beta(1)\eta_{\mathrm{alt}}(w-1).
\end{align*}
Substitution in \eqref{newmixed:eq:two-d} and division by two prove
\eqref{newmixed:eq:family}.  The displayed examples use
$\beta(1)=\pi/4$, $\beta(3)=\pi^3/32$, and
$\beta(5)=5\pi^5/1536$.
\end{proof}

For an independent numerical check, the defining iterated integral gives
\begin{equation}\label{newmixed:eq:integral}
 u_{w-1,1}=\Im\left[
 \frac1{(w-2)!}\int_0^1
 \frac{-i\log(1+it)}{1-it}(-\log t)^{w-2}\,dt\right].
\end{equation}
The path meets no singularity.  Expansion in the open disk, followed
by the ordinary boundary limit, verifies this representation.  Its
integrand is $O(t|\log t|^{w-2})$ at zero, so it is integrable even
at $w=2$.  Independent quadrature of \eqref{newmixed:eq:integral}
agrees with \eqref{newmixed:eq:family} at weights $2,4,6,8,10$ to the
working precision of seventy decimal digits.  In particular
\[
 u_{3,1}\approx0.01508330704224296705,
 \qquad u_{5,1}\approx0.001899710391767247024.
\]
The proof is the finite polynomial calculation above; the quadrature is
a diagnostic.  Classical colored parity theorems predict reducibility
in this sector.  The formula supplies a compact explicit family and
a direct proof; no claim of worldwide priority is required for that
consequence.

\paragraph{A short integer row certificate.}
Using the left-side row notation of
\eqref{newrank:eq:stuffle}--\eqref{newrank:eq:shuffle}, the entire family
has the explicit coefficient certificate
\begin{align}
 2x_{w-1,1}^{1,2}
 ={}&\sum_{p=1}^{w-1}\mathsf S_{p,w-p}^{2,1}
   +\mathsf H_{1,w-1}^{2,1}
   +\sum_{p=1}^{w-1}(-1)^{w-p-1}\mathsf H_{p,w-p}^{2,1}
   +\mathsf S_{1,w-1}^{2,1}\notag\\
 &-\mathsf S_{w-1,1}^{1,1}-\mathsf H_{1,w-1}^{1,3}.
 \label{newmixed:eq:rowcertificate}
\end{align}
All rows on the right come from finite products.  This identity is
understood after the manuscript's color-conjugation identifications.
Its coefficients are integers in $\{-1,1\}$ before equal rows are
collected, and it uses $2w+2$ listed row occurrences.  The companion
script checks the exact coefficient cancellation through weight twenty.
